% PAGES: 183-184
% Continues directly after the "Compute the first row of U:" bullet opened
% at the end of sec06_c1_3.tex (page 182 / folio 166); that itemize was
% already closed there, so the pseudocode box below is plain text, not a
% list item.
%
% THIS FILE ENDS WITH AN OPEN MANUAL EXAMPLE BLOCK, still inside a
% "\begin{examplex}" opened below on page 184 / folio 168 ("Example: Let A
% be the following symmetric matrix:"). The example computes a full 4x4
% Cholesky decomposition and is NOT finished on folio 168 -- the last line
% is an unnumbered display equation ending in a trailing comma
% ("(U(2:4,2:4))^tU(2:4,2:4) = [matrix],") that continues onto the next
% printed page (folio 169, PDF page 185), which is outside this chunk. The
% next chunk agent must NOT open a new "\begin{examplex}" or reprint the
% "Example:" label -- it should continue the sentence immediately after
% this file's final equation and supply the environment's auto-appended
% closing square only once the source's own closing square is reached.
% scripts/check_env_balance.py WILL flag this "examplex" as unclosed at
% chapter end until the next chunk supplies "\end{examplex}".

\begin{center}
\fbox{%
\begin{minipage}{0.5\linewidth}
$U(1,1) = \sqrt{A(1,1)}$ \\
for $k = 2:n$ \\
\hspace*{1.5em} $U(1,k) \;=\; \dfrac{A(1,k)}{U(1,1)}$ \\
end
\end{minipage}%
}
\end{center}

\begin{itemize}
\item Update the $(n-1) \times (n-1)$ lower right part of $A$ as follows:
\end{itemize}
\begin{equation}
A(2:n,2:n) \;=\; A(2:n,2:n) \;-\; (U(1,2:n))\tp U(1,2:n),
\end{equation}
see (6.21), which can be written entry by entry as follows:

\begin{center}
\fbox{%
\begin{minipage}{0.6\linewidth}
for $j = 2:n$ \\
\hspace*{1.5em} for $k = j:n$ \\
\hspace*{3em} $A(j,k) = A(j,k) - U(1,j)U(1,k)$ \\
\hspace*{1.5em} end \\
end
\end{minipage}%
}
\end{center}

Note that in the second ``for'' loop above, only the upper triangular part of the
matrix is updated, i.e., we use ``for $k = j:n$'' instead of ``for $k = 2:n$'', as was
the case for the LU decomposition without pivoting, see the pseudocode from Table 2.5,
since the matrix $A(2:n,2:n) - (U(1,2:n))\tp U(1,2:n)$ is symmetric. This accounts for
the computational savings in the Cholesky decomposition compared to the LU
decomposition; see Lemma 6.2 and Lemma 2.3.

Every row of $U$ is thereafter computed recursively from the latest updated part of the
matrix $A$; see also the $4 \times 4$ example below. For example, to compute the
$i$-th row of $U$, we do the following:

\begin{itemize}
\item Compute the $i$-th row of $U$:
\end{itemize}

\begin{center}
\fbox{%
\begin{minipage}{0.55\linewidth}
$U(i,i) = \sqrt{A(i,i)}$ \\
for $k = (i+1):n$ \\
\hspace*{1.5em} $U(i,k) \;=\; \dfrac{A(i,k)}{U(i,i)}$ \\
end
\end{minipage}%
}
\end{center}

\begin{itemize}
\item Update the $(n-i) \times (n-i)$ lower right part of $A$ as follows:
\end{itemize}
\begin{equation}
\begin{split}
A(i+1:n,i+1:n) \;&=\; A(i+1:n,i+1:n) \\
&\quad -\; (U(i,i+1:n))\tp U(i,i+1:n),
\end{split}
\end{equation}
which can be written entry by entry as follows:

\begin{center}
\fbox{%
\begin{minipage}{0.65\linewidth}
for $j = (i+1):n$ \\
\hspace*{1.5em} for $k = j:n$ \\
\hspace*{3em} $A(j,k) = A(j,k) - U(i,j)U(i,k)$; \\
\hspace*{1.5em} end \\
end
\end{minipage}%
}
\end{center}

In the second ``for'' loop above, only the upper triangular part of the matrix is
updated, i.e., we use ``for $k = j:n$'' instead of ``for $k = (i+1):n$'', since the
matrix $A(i+1:n,i+1:n) - (U(i,i+1:n))\tp U(i,i+1:n)$ is symmetric.

Further clarification on the recursive part of the Cholesky decomposition algorithm
can be found in the example below for a $4 \times 4$ matrix.

\begin{examplex}
Let $A$ be the following symmetric matrix:
\begin{equation}
A \;=\; \begin{pmatrix} 9 & -3 & 6 & -3 \\ -3 & 5 & -4 & 7 \\ 6 & -4 & 21 & 3 \\ -3 & 7 &
3 & 15 \end{pmatrix}.
\end{equation}
We attempt to find the Cholesky factor $U$ of the matrix $A$ using the recursive
method described above. If this method fails, then $A$ is not a symmetric positive
definite matrix. If the method succeeds, the Cholesky factor of $A$ is found.

The first row of $U$ is computed as follows: $U(1,1) = \sqrt{A(1,1)} = 3$; cf. (6.6).
From (6.8), we find that $U(1,k) = A(1,k)/U(1,1)$ for all $k = 2:4$, and therefore
\begin{align*}
U(1,2) = \frac{-3}{U(1,1)} = \frac{-3}{3} = -1; \quad U(1,3) &= \frac{6}{U(1,1)} =
\frac{6}{3} = 2; \\
U(1,4) = \frac{-3}{U(1,1)} = \frac{-3}{3} &= -1.
\end{align*}
Then, the current form of $U$ is
\begin{equation}
U \;=\; \begin{pmatrix} 3 & -1 & 2 & -1 \\ 0 & U(2,2) & U(2,3) & U(2,4) \\ 0 & 0 &
U(3,3) & U(3,4) \\ 0 & 0 & 0 & U(4,4) \end{pmatrix}.
\end{equation}
The updated form of the $3 \times 3$ matrix $A(2:4,2:4)$ is computed using (6.28), from
$A(2:4,2:4)$ obtained from (6.30) and with $U$ given by (6.31), as follows:
\begin{align*}
A(2:4,2:4)
&= A(2:4,2:4) \;-\; (U(1,2:4))\tp U(1,2:4) \\
&= \begin{pmatrix} 5 & -4 & 7 \\ -4 & 21 & 3 \\ 7 & 3 & 15 \end{pmatrix} \;-\;
\begin{pmatrix} -1 \\ 2 \\ -1 \end{pmatrix} \begin{pmatrix} -1 & 2 & -1 \end{pmatrix} \\
&= \begin{pmatrix} 5 & -4 & 7 \\ -4 & 21 & 3 \\ 7 & 3 & 15 \end{pmatrix} \;-\;
\begin{pmatrix} 1 & -2 & 1 \\ -2 & 4 & -2 \\ 1 & -2 & 1 \end{pmatrix} \\
&= \begin{pmatrix} 4 & -2 & 6 \\ -2 & 17 & 5 \\ 6 & 5 & 14 \end{pmatrix}.
\end{align*}
Thus, the updated form of the $3 \times 3$ matrix $A(2:4,2:4)$ is
\begin{equation}
A(2:4,2:4) \;=\; \begin{pmatrix} 4 & -2 & 6 \\ -2 & 17 & 5 \\ 6 & 5 & 14
\end{pmatrix}.
\end{equation}
Then,
\[
(U(2:4,2:4))\tp U(2:4,2:4) \;=\; \begin{pmatrix} 4 & -2 & 6 \\ -2 & 17 & 5 \\ 6 & 5 & 14
\end{pmatrix},
